\section{The sharp homogeneous angular coefficient}\label{sec:angular}
The earlier Fourier proof remains useful when the external test law is
required to have independent single commands. Its elementary angular
estimate admits the following sharpened statement. This also resolves
the coefficient issue in the preceding report without changing any
previous published source.

\begin{lemma}[Homogeneous angular averaging]\label{lem:angular}
Let $f_l(re^{i\theta})=r e^{il\theta}$ for $r>0$ and $f_l(0)=0$.
For every integrable complex $Z$ and integer $l\ge1$,
\begin{equation}\label{eq:angular-sharp}
 |\E f_l(Z)-f_l(\E Z)|
       \le(l^2-1)(\E|Z|-|\E Z|).
\end{equation}
For $l\ge2$ the coefficient $l^2-1$ cannot be decreased uniformly over
all such mixtures. At $l=1$ both sides vanish.
\end{lemma}
\begin{proof}
The case $l=1$ is linear. For $l\ge2$, if $\E Z=0$, the left side is
at most $\E|Z|$ and $l^2-1\ge1$. Otherwise rotate the variable so
$\E Z=\rho>0$, and write $Z=re^{i\theta}$. For $u=e^{i\theta}$,
\[
 u^l-1-l(u-1)=(u-1)\sum_{j=0}^{l-1}(u^j-1).
\]
Since $|u^j-1|\le j|u-1|$, the modulus of this remainder is at most
$l(l-1)(1-\cos\theta)$. Retaining this coefficient gives
\[
 |u^l-\cos\theta-il\sin\theta|
       \le\{l(l-1)+(l-1)\}(1-\cos\theta)
       =(l^2-1)(1-\cos\theta).
\]
Multiply by $r$ and average. The identities
$\E(r\sin\theta)=0$ and $\E(r\cos\theta)=\rho$ prove
\eqref{eq:angular-sharp}. Rotation covariance restores the general case.
For sharpness take $Z=e^{i\theta}$ and $e^{-i\theta}$ with equal
probabilities, and let $\theta\downarrow0$. The ratio is
\[
       \frac{|\cos(l\theta)-\cos\theta|}{1-\cos\theta}
                          \longrightarrow l^2-1.
\]
\end{proof}
Thus a uniformly linear-in-$l$ replacement is impossible for unrestricted
conditional mixtures. The improvement to cube-root width in this paper
does not require such an estimate: it changes the test law to packets
and measures the directly incurred geometric amplitude loss.
